\documentclass[10pt,aspectratio=169]{beamer}
\usetheme{metropolis}
\usepackage{booktabs}
\usepackage{tikz}
\usetikzlibrary{positioning,arrows.meta,shapes,fit,backgrounds,calc}
\usepackage{amsmath,amssymb}
\usepackage{xcolor}

\definecolor{accent}{HTML}{0066CC}
\definecolor{accent2}{HTML}{CC3300}
\definecolor{accent3}{HTML}{009933}
\definecolor{lightgray}{HTML}{F0F0F0}

\title{GRU Prefetcher V10: From Accuracy Sweep to Utility Sweep}
\subtitle{Reuse distance, cache state, phase, and bypass gating}
\author{Angela Wu}
\institute{CPU Cache + NN Prefetching Research}
\date{Update \#3}

\begin{document}

% =========================================================
\begin{frame}
  \titlepage
\end{frame}

% =========================================================
\begin{frame}{Are We on the Right Track?}
  \textbf{Yes -- but the target has changed.}

  \vspace{0.4em}
  The right research method is:
  \begin{enumerate}
    \item Add one input feature at a time: PC, PC$\Delta$, reuse distance, cache pressure, phase.
    \item Record both ML metrics and system metrics: accuracy/F1, CPU inference latency, trigger rate, IPC, MPKI, bandwidth/pollution.
    \item Use the sweep itself to discover which feature matters on which workload.
  \end{enumerate}

  \vspace{0.5em}
  \textbf{The correction:} accuracy is not the final goal.
  The hardware question is not ``can the GRU predict the next address?''
  It is:
  \[
      \text{Should this candidate prefetch / insert / bypass action be taken under current cache resources?}
  \]

  \vspace{0.4em}
  \textbf{Therefore V10 changes the formulation from address prediction to utility-controlled candidate prediction.}
\end{frame}

% =========================================================
\begin{frame}{What the Current Data Already Told Us}
  \textbf{Main insight: better ML accuracy can still reduce IPC.}

  \vspace{0.3em}
  From V8/V9:
  \begin{itemize}
    \item V8 improved top-1 accuracy dramatically, but IPC was still below baseline on the held-out replay.
    \item V9 moved to in-distribution train/val/test and delta-bitmap output; this fixed the cross-binary mismatch, but not the system-level question.
    \item The important number is no longer just top-1/F1. The important pair is:
      \[
        (\text{prediction quality},\ \text{prefetch usefulness under cache/bandwidth/MSHR constraints})
      \]
  \end{itemize}

  \vspace{0.4em}
  \textbf{Interpretation:}
  \begin{itemize}
    \item If accuracy rises but IPC falls, the model is issuing too many wrong/late/early/polluting prefetches.
    \item If top-1 is high but trigger rate is too high, the model is acting like an aggressive prefetcher.
    \item If SPP beats GRU on lbm/gcc, the missing signal is probably not ``more GRU layers''; it is cache-state / reuse / timeliness.
  \end{itemize}
\end{frame}

% =========================================================
\begin{frame}{Paper-Grounded Design Rule}
  \textbf{The papers point to the same conclusion from different directions.}

  \vspace{0.3em}
  \begin{center}
  \footnotesize
  \begin{tabular}{p{3.0cm}p{4.7cm}p{5.3cm}}
    \toprule
    \textbf{Paper line} & \textbf{What it says} & \textbf{What V10 uses} \\
    \midrule
    Hashemi / Voyager & Neural prefetching should be classification, not address regression; page/offset helps reduce class explosion. & Keep bounded outputs: delta bitmap / candidate score. \\
    Twilight & Do not directly predict arbitrary addresses; select among likely candidates and include no-prefetch. & Add explicit confidence gate and no-prefetch behavior. \\
    APT-GET & Timeliness matters: a correct prefetch can still be useless if too early/late. & Add reuse-distance and phase/timeliness features. \\
    Limoncello & Prefetch can hurt under bandwidth pressure. & Track trigger rate and add bypass/suppress gate. \\
    DART / Net2Tab & Heavy NN is for offline learning; hardware path should become table-like. & GRU is an exploration tool; final design should be counters/tables/small scorer. \\
    Alecto / page-size work & Best policy changes by workload and phase. & Sweep traces and phase features, not one universal setting. \\
    \bottomrule
  \end{tabular}
  \end{center}
\end{frame}

% =========================================================
\begin{frame}{Why GRU Is Still the Right First NN}
  \textbf{Use GRU as the search model, not the final hardware block.}

  \vspace{0.3em}
  D2L's GRU has two gates:
  \[
  R_t = \sigma(X_t W_{xr} + H_{t-1}W_{hr}+b_r),\qquad
  Z_t = \sigma(X_t W_{xz} + H_{t-1}W_{hz}+b_z)
  \]
  \[
  \tilde{H}_t = \tanh(X_t W_{xh} + (R_t \odot H_{t-1})W_{hh}+b_h)
  \]
  \[
  H_t = Z_t\odot H_{t-1} + (1-Z_t)\odot\tilde{H}_t
  \]

  \vspace{0.3em}
  \textbf{Hardware interpretation:}
  \begin{itemize}
    \item Reset gate $\approx$ when the recent delta history should be forgotten.
    \item Update gate $\approx$ when the long-term per-PC behavior should persist.
    \item This matches memory streams: some PCs have stable stride, some reset at phase boundaries.
  \end{itemize}

  \vspace{0.3em}
  \textbf{But:} online CPU inference must eventually be compressed to lookup tables, counters, low-bit MLP, or a tabularized scorer.
\end{frame}

% =========================================================
\begin{frame}{V10 Model: Behavior-Aware GRU + Bypass Gate}
  \begin{center}
  \begin{tikzpicture}[
      node distance=0.5cm,
      box/.style={draw,rounded corners,minimum width=2.5cm,minimum height=0.75cm,font=\scriptsize,align=center},
      data/.style={fill=lightgray},
      model/.style={fill=accent!15},
      out/.style={fill=accent3!18},
      bad/.style={fill=accent2!15}]
    \node[box,data] (hist) {delta history\\last $H$ tokens};
    \node[box,data,right=of hist] (pc) {PC / PC$\Delta$\\hash};
    \node[box,data,right=of pc] (reuse) {reuse distance\\bucket};
    \node[box,data,right=of reuse] (state) {cache state\\phase bucket};

    \node[box,model,below=0.8cm of pc] (gru) {small GRU\\hidden 64};
    \node[box,model,below=0.8cm of reuse] (trunk) {tiny MLP scorer};

    \node[box,out,below=0.8cm of trunk,xshift=-1.8cm] (bm) {delta bitmap\\candidate scores};
    \node[box,bad,below=0.8cm of trunk,xshift=1.8cm] (bp) {bypass / suppress\\gate};

    \draw[->] (hist) -- (gru);
    \draw[->] (pc) -- (trunk);
    \draw[->] (reuse) -- (trunk);
    \draw[->] (state) -- (trunk);
    \draw[->] (gru) -- (trunk);
    \draw[->] (trunk) -- (bm);
    \draw[->] (trunk) -- (bp);
  \end{tikzpicture}
  \end{center}

  \vspace{0.3em}
  \textbf{Decision rule at export time:}
  \[
    \text{emit}(a+\Delta) =
    [p_{\Delta} \ge \theta_{\Delta}]
    \land
    [p_{\text{bypass/suppress}} < \theta_b]
    \land
    [\text{degree} \le D_{\max}]
  \]
\end{frame}

% =========================================================
\begin{frame}{Controlled Feature Sweep Plan}
  \textbf{Each row changes exactly one thing.}

  \vspace{0.4em}
  \begin{center}
  \scriptsize
  \begin{tabular}{llll}
    \toprule
    \textbf{Variant} & \textbf{Input features} & \textbf{Output} & \textbf{Question answered} \\
    \midrule
    V10a & delta history only & bitmap & Can sequence alone learn useful deltas? \\
    V10b & + PC bucket & bitmap & Does instruction identity help in-distribution? \\
    V10c & + PC$\Delta$ bucket & bitmap & Does local context improve SPP-like behavior? \\
    V10d & + reuse-distance bucket & bitmap & Does cacheline behavior predict usefulness? \\
    V10e & + cache pressure + phase & bitmap & Does system state reduce harmful prefetches? \\
    V10f & + bypass/suppress head & bitmap + gate & Can we keep accuracy while reducing pollution? \\
    \bottomrule
  \end{tabular}
  \end{center}

  \vspace{0.5em}
  \textbf{Metric table to fill for every trace:}
  \[
  \text{variant},\ \text{params},\ \text{F1},\ \text{top1/top5},\
  \text{CPU $\mu$s},\ \text{trigger rate},\ \text{avg degree},\
  \text{IPC},\ \text{MPKI},\ \text{speedup}
  \]
\end{frame}

% =========================================================
\begin{frame}{Recommended GRU Parameters}
  \begin{center}
  \footnotesize
  \begin{tabular}{lll}
    \toprule
    \textbf{Knob} & \textbf{V10 default} & \textbf{Why} \\
    \midrule
    History length & 16 & Long enough for irregular deltas; still cheap. \\
    Delta vocab & 2048 + OOV & V8 top-256 had too much OOV; larger vocab reduces collapse. \\
    Delta output & bitmap over $[-64,+64]$ & Bounded candidate set; can issue 0--2 prefetches. \\
    GRU hidden & 64 first, sweep 32/64/96 & 64 is the hardware-friendly default. \\
    Delta embedding & 24--32 & Enough to separate common delta classes. \\
    PC buckets & 1024 & Useful in-distribution; not trusted cross-binary. \\
    PC$\Delta$ buckets & 4096 & Captures local context similar to signature/path predictors. \\
    Reuse buckets & 16 log buckets & Tiny hardware counter feature. \\
    Cache/phase buckets & 8--16 & Enough for pressure/phase without overfitting. \\
    Thresholds & sweep $\theta_\Delta,\theta_b$ & IPC depends on aggressiveness, not only F1. \\
    \bottomrule
  \end{tabular}
  \end{center}

  \vspace{0.4em}
  \textbf{Do not increase layers first.}
  If V10 underperforms SPP, add better hardware features / gating before making the GRU bigger.
\end{frame}

% =========================================================
\begin{frame}{Reuse Distance Feature}
  \textbf{Definition:} for current cache line $\ell_i$,
  \[
      RD_i =
      \begin{cases}
      i - \mathrm{last}(\ell_i), & \text{if line was seen before}\\
      \infty, & \text{otherwise}
      \end{cases}
  \]

  \vspace{0.3em}
  Use log buckets:
  \[
      b_i = \min(B-1,\ \lfloor \log_2(RD_i + 1) \rfloor)
  \]

  \vspace{0.3em}
  \textbf{Why this matters:}
  \begin{itemize}
    \item Small reuse distance $\Rightarrow$ hot line / likely useful / keep or prefetch.
    \item Huge or infinite reuse distance $\Rightarrow$ streaming or dead line / bypass or low-priority insert.
    \item This directly connects the NN feature to cache replacement and bypass, not just prefetch address prediction.
  \end{itemize}

  \vspace{0.3em}
  \textbf{Hardware version:} a small sampled line table with saturating/log counters, not a full exact reuse-distance tracker.
\end{frame}

% =========================================================
\begin{frame}{Trace Choice}
  \textbf{The current 3 traces are good for first sweep.}

  \vspace{0.3em}
  \begin{center}
  \footnotesize
  \begin{tabular}{lll}
    \toprule
    \textbf{Trace} & \textbf{Why keep it} & \textbf{What it should reveal} \\
    \midrule
    \texttt{605.mcf\_s-994B} & irregular / pointer-heavy & Whether reuse/bypass beats raw PC. \\
    \texttt{619.lbm\_s-4268B} & streaming / stencil-like & Whether bitmap + degree control can approach SPP. \\
    \texttt{602.gcc\_s-734B} & mixed control + memory behavior & Whether phase/cache-state features matter. \\
    \bottomrule
  \end{tabular}
  \end{center}

  \vspace{0.4em}
  \textbf{Add next if time allows:}
  \begin{itemize}
    \item \texttt{620.omnetpp\_s}: object/pointer-heavy, good for overfitting test.
    \item one more high-SPP-headroom trace from slide 8: choose the trace where SPP gives the biggest gap over LRU.
  \end{itemize}

  \vspace{0.3em}
  \textbf{Do not judge novelty from one trace.}
  The idea is strongest if V10 explains \emph{why} different features win on different workload classes.
\end{frame}

% =========================================================
\begin{frame}{Creative Idea Emerging From the Sweep}
  \textbf{Behavior-Class Utility Predictor (BCUP), not just GRU prefetcher.}

  \vspace{0.3em}
  The GRU sweep is a discovery tool. The final research contribution can be:

  \vspace{0.3em}
  \begin{center}
  \Large
  \textbf{A tiny hardware utility controller that learns when to prefetch, bypass, or suppress candidates based on behavior class.}
  \end{center}

  \vspace{0.4em}
  Candidate generators can stay conventional:
  \[
    \text{next-line},\ \text{stride},\ \text{SPP},\ \text{BO},\ \text{temporal/page successor}
  \]

  \vspace{0.3em}
  The learned part only scores:
  \[
    U(c \mid \text{PC}, \Delta\text{-history}, RD, \text{phase}, \text{cache pressure})
  \]

  \vspace{0.3em}
  This avoids the hard problem of predicting arbitrary addresses and targets the hardware problem directly: \textbf{utility under constraints}.
\end{frame}

% =========================================================
\begin{frame}{Immediate Implementation Status}
  \textbf{New notebook:} \texttt{notebook/gru\_sweep\_v10\_reuse\_phase.ipynb}

  \vspace{0.3em}
  It adds:
  \begin{itemize}
    \item V10a--V10f controlled feature variants.
    \item Reuse-distance buckets.
    \item Cache-pressure / phase buckets.
    \item Optional bypass/suppress head.
    \item Summary CSV for accuracy/F1/latency/trigger-rate.
    \item Prefetch-list export compatible with the existing replayer format:
      \[
        \texttt{idx\ \ 0xprefetch\_addr}
      \]
  \end{itemize}

  \vspace{0.3em}
  \textbf{Run path after Colab export:}
  \[
    \texttt{PFETCH=results/generated/prefetch\_lists/<file>.txt\ MODEL\_TAG=GRU\_V10... bash scripts/run\_nn\_replay.sh}
  \]

  \vspace{0.3em}
  The \texttt{.sh} files are left unchanged.
\end{frame}

\end{document}
